% Q1 Part 2 and appendix -- CE experiment (roundabout, frozen DT + LLM supervisor)
% Red (\rev / revblock) = changed or new with respect to the previous version of this experiment (prompt v5);
% black = unchanged (platform, DT, executor, metric definitions).
% Model inputs and responses are verbatim from results/hazard_test_v6/decisions_llm_llm5.csv; long lines are wrapped
% (continuation lines indented by 4 spaces), nothing else is changed.
\documentclass[11pt]{article}
\usepackage[margin=1in]{geometry}
\usepackage[T1]{fontenc}
\usepackage{lmodern}
\usepackage{xcolor}
\usepackage{booktabs}
\usepackage{fancyvrb}
\usepackage{array}
\usepackage{tabularx}
\usepackage{enumitem}
\usepackage{hyperref}
\newcommand{\rev}[1]{\textcolor{red}{#1}}
\newenvironment{revblock}{\color{red}}{}
% one style for all model inputs and responses: same font size and indentation, distinguished by the frame label
\DefineVerbatimEnvironment{ModelInput}{Verbatim}{fontsize=\footnotesize, xleftmargin=1.5em, frame=single,
  framesep=2mm, label={\textbf{Model input (user message, verbatim)}}}
\DefineVerbatimEnvironment{ModelResponse}{Verbatim}{fontsize=\footnotesize, xleftmargin=1.5em, frame=lines,
  framerule=0.8pt, framesep=2mm, label={\textbf{Model response (raw output, verbatim)}}}
\DefineVerbatimEnvironment{PromptText}{Verbatim}{fontsize=\footnotesize, xleftmargin=1.5em, frame=single,
  framesep=2mm, label={\textbf{System prompt (verbatim)}}}

\begin{document}

\section*{Q1 Part 2: Implementation and Evaluation on the Roundabout Scenario}

\subsection*{2.1 \rev{Hypothesis and why the previous evidence was not sufficient}}

\begin{revblock}
\paragraph{Hypothesis.} A Decision Transformer (DT) trained only on demonstrations of ordinary traffic does not cover
situations outside its training distribution. A pretrained LLM may already carry the relevant scene understanding,
traffic common sense and reasoning. We test whether this capability, used at run time as a high-level supervisor of a
\emph{fixed} DT, lets the DT handle out-of-distribution (OOD) situations without additional policy training and
without degrading its nominal performance. The supervisor has to judge what a situation means for the ego's own route,
what the DT's proposed action will lead to, whether to intervene, with which maneuver, and on which vehicle.

\paragraph{What the previous version (prompt v5) could not show.} Checked against the code, the previous version
(i)~told the model which vehicles were ``in your path'' and labelled them \texttt{STATIONARY} or ``entering the
single-lane exit ahead of you'', (ii)~gave a rule table with examples that matched the test conditions one to one,
and (iii)~let the executor, not the model, choose the target (the slowest of the first three path vehicles). In
addition, every abnormal vehicle in the test set was on the ego's path, so ``intervene whenever something abnormal is
listed'' was always correct. Its results therefore show that the integrated system works, but they cannot be
attributed to the model's own understanding.

\paragraph{Minimal changes.} (1)~Two control scenarios with the same stalled car placed where it does \emph{not}
affect the ego; (2)~a facts-only input without labels; (3)~a prompt that defines the task but contains no rule table
and no examples; (4)~the model returns the target vehicle id; (5)~mechanism controls and a paired evaluation on new,
untouched seeds. The DT stays frozen and the executor is unchanged.
\end{revblock}

\subsection*{2.2 Why the DT needs scene understanding}

The ego is driven by a DT trained on 5409 trajectories of the unmodified \texttt{roundabout-v0} environment
(highway-env, occupancy-grid observation, five discrete meta-actions at 2\,Hz, 11\,s episodes), in which all traffic
flows at about 16\,m/s. The checkpoint is frozen and the target return is the one used in training.
\begin{revblock}
A stopped or crawling car never occurs in these data. When such a car is on the ego's route, the DT collides in
100\% of the test episodes (Table~\ref{tab:closed}). The same stopped car placed after the ego's exit or upstream of
its merge point does not matter: the DT completes 88\% and 100\% of those episodes. Detecting an abnormal object is
therefore not enough. Whether and how to intervene depends on the object's position relative to the ego's route, its
motion, the DT's proposed action and the road topology (a single-lane exit cannot be entered in front of a car that
is already entering it). This is the judgment we expect from the LLM.
\end{revblock}

\subsection*{2.3 \rev{What the model has to judge}}

\begin{revblock}
At every policy step (0.5\,s) the model must decide (a)~which vehicles, if any, can affect the part of the road the ego
will actually drive, (b)~what happens if the DT's proposed action is executed, (c)~whether to keep the DT
(\texttt{FOLLOW\_POLICY}), pass a vehicle in the other ring lane (\texttt{BYPASS}) or stay behind it without changing
lanes (\texttt{FOLLOW\_VEHICLE}), and (d)~which vehicle the maneuver refers to. The scenarios are built as minimal
pairs, each differing in one relation only:
\begin{itemize}[nosep]
\item before vs.\ after the exit: stalled car 9.6\,m before (\texttt{stall\_exit}, bypass) or 9.6\,m after the ego's exit
  in the same lane (\texttt{stall\_past\_exit}, no intervention);
\item stationary vs.\ crawling at the same place: 0\,m/s (\texttt{stall\_exit}, bypass) vs.\ 3\,m/s, about to enter the
  single-lane exit (\texttt{slow\_exit}, follow, it cannot be overtaken);
\item ahead vs.\ behind: stalled car 27.6\,m after the merge point (\texttt{stall\_mid}, bypass) vs.\ 12\,m upstream of
  it (\texttt{stall\_upstream}, no intervention).
\end{itemize}
Nominal traffic is kept as the control for unnecessary interventions.
\end{revblock}

\subsection*{2.4 \rev{Design of the OOD validation}}

\paragraph{Architecture.} The frozen DT proposes an action at every step. The supervisor either keeps it or starts a
temporally extended maneuver; a maneuver executor turns the maneuver into actions (rate-limited braking with a
safe-distance law, lane change only with a speed-dependent gap in the target lane, end of the maneuver once the target
is 8\,m behind the ego). \rev{The executor only checks that the returned target id exists in the current perception
list; otherwise it refuses, keeps the DT action and logs the refusal. It no longer selects the target and makes no
relevance judgment: a maneuver on an irrelevant vehicle is executed and its cost appears in the results.}

\begin{revblock}
\paragraph{Input (facts only).} Ego road, lane and speed; navigation (where the north exit leaves the ring); the DT's
proposed action and current target speed; and every vehicle within 50\,m, grouped by road and lane (lane association
is a perception output), with a stable id, its position along the ring (signed, in the driving direction), its speed
and, for ring vehicles, the time the ego needs at current speeds to reach its position (computed identically for every
lane). Map landmarks (\texttt{[YOU]}, \texttt{[YOUR MERGE POINT]}, \texttt{[YOUR EXIT]}) are placed in the same
coordinate. The input contains no motion-state labels, no ``in your path'' grouping, no exit-zone or predicted-maneuver
labels, no hazard marker and no scenario name. Since this is a text LLM, the task is scene judgment from structured
observations, not perception from images.

\paragraph{Prompt (task only).} Road facts, the DT's training limitation, the meaning and execution constraints of the
three decisions, the statement that every intervention has a cost, and the task: judge the consequence of the DT's
action for the ego's own route, then decide. No rule table and no examples (full text in Appendix B).
Output: \texttt{\{"consequence", "decision", "target\_id"\}}; the model is not asked to repeat its input.

\paragraph{Controls (same facts, same executor).}
\begin{itemize}[nosep]
\item \textbf{DT}: baseline.
\item \textbf{Naive trigger}: any vehicle slower than 5\,m/s within 40\,m in front of the ego (straight-line geometry,
  no lane or route reasoning) $\rightarrow$ \texttt{BYPASS}. Measures what a hand-written avoidance logic alone gives.
\item \textbf{LLM v5}: the previous pipeline unchanged (labels, rule prompt, executor-chosen target). Measures what a
  prompt and input that already contain the answer give.
\item \textbf{LLM v6}: facts-only input, task-only prompt, model-chosen target (the method under test).
\item \textbf{Designed rule}: the designer's route-aware rule (the answer key). Upper reference for correct decisions
  with this executor, not a competitor.
\end{itemize}

\paragraph{Protocol.} Development on seeds 200--219 only (offline single-step checks and one closed-loop pipeline
check). The input, prompt, model (Qwen3-4B-Instruct-2507, Q4\_K\_M, temperature 0, Ollama on an RTX~3070~Ti) and
executor were frozen and the protocol was written down before any test run. Test: seeds 600--649, 50 episodes per
scenario, identical scenarios and seeds for all methods (paired), 11\,s horizon. Outcomes: success, collision,
at-fault collision (the ego is not at fault only if rear-ended while keeping its lane outside the merge zone),
timeout, speed, completion time. Mechanism measures: every supervisor call is compared with a privileged reference that
is never shown to any supervisor (intervene on the injected vehicle while it is in the ego's lane ahead and before the
exit); maneuvers started on the injected vehicle; unnecessary maneuvers. Paired comparisons use the exact McNemar test.
A scoring-code bug found after freezing (no effect on any driving decision) was fixed and all methods were rerun; the
rule controls were identical episode by episode before and after.
\end{revblock}

\subsection*{2.5 \rev{Results}}

\begin{revblock}
\begin{table}[h]
\centering\footnotesize
\caption{Closed loop on test seeds 600--649 (50 episodes per cell): success / collision / at-fault collision / timeout (\%).
``Need'': stall\_mid, stall\_exit, slow\_exit; ``no need'': stall\_past\_exit, stall\_upstream.}
\label{tab:closed}
\setlength{\tabcolsep}{3pt}
\begin{tabular}{lcccccc}
\toprule
Method & nominal & stall\_mid & stall\_exit & slow\_exit & stall\_past\_exit & stall\_upstream \\
\midrule
DT & 94/6/6/0 & 0/100/100/0 & 0/100/100/0 & 0/100/100/0 & 88/12/12/0 & 100/0/0/0 \\
Naive trigger & 94/6/6/0 & 66/20/20/14 & 70/6/4/24 & 0/80/80/20 & 80/6/6/14 & 100/0/0/0 \\
LLM v5 (labels + rules) & 94/6/6/0 & 66/20/20/14 & 58/4/4/38 & 0/10/6/90 & 64/6/6/30 & 100/0/0/0 \\
\textbf{LLM v6 (facts only)} & 96/4/4/0 & 62/24/24/14 & 54/34/34/12 & 0/80/80/20 & 84/4/4/12 & 100/0/0/0 \\
Designed rule (answer key) & 94/6/6/0 & 66/20/20/14 & 68/6/4/26 & 0/18/12/82 & 88/12/12/0 & 100/0/0/0 \\
\bottomrule
\end{tabular}

\vspace{4pt}
\begin{tabular}{lcc}
\toprule
Method & mean, need & mean, no need \\
\midrule
DT & 0 / 100 / 100 / 0 & 94 / 6 / 6 / 0 \\
Naive trigger & 45.3 / 35.3 / 34.7 / 19.3 & 90 / 3 / 3 / 7 \\
LLM v5 (labels + rules) & 41.3 / 11.3 / 10.0 / 47.3 & 82 / 3 / 3 / 15 \\
\textbf{LLM v6 (facts only)} & 38.7 / 46.0 / 46.0 / 15.3 & 92 / 2 / 2 / 6 \\
Designed rule (answer key) & 44.7 / 14.7 / 12.0 / 40.7 & 94 / 6 / 6 / 0 \\
\bottomrule
\end{tabular}
\end{table}

No method can succeed in slow\_exit within 11\,s (following a 3\,m/s car to the exit takes longer), so there only
collisions separate the methods.

\begin{table}[h]
\centering\footnotesize
\caption{Decision quality of LLM v6 against the privileged reference, and the first maneuver started on the injected
vehicle (episodes out of 50; BYPASS/FOLLOW in brackets).}
\label{tab:dec}
\setlength{\tabcolsep}{4pt}
\begin{tabular}{>{\raggedright\arraybackslash}p{4.6cm}cccccc}
\toprule
 & stall\_mid & stall\_exit & slow\_exit & stall\_past\_exit & stall\_upstream & nominal \\
\midrule
v6: correct decision and target, when the reference intervenes & 60\% & 8\% & 0\% & -- & -- & -- \\
v6: intervenes when the reference does not & 26\% & 52\% & 49\% & 48\% & 48\% & 35\% \\
v6: episodes with $\geq$1 maneuver & 50 & 50 & 50 & 50 & 50 & 50 \\
v6: first maneuver on the injected car & 36 (36/0) & 10 (10/0) & 3 (3/0) & 19 & 15 & 0 \\
v5: first maneuver on the injected car & 50 (50/0) & 48 (38/10) & 46 (1/45) & 48 & 0 & 0 \\
Designed: first maneuver on the injected car & 50 (50/0) & 49 (49/0) & 44 (0/44) & 3 & 0 & 0 \\
\bottomrule
\end{tabular}
\end{table}

LLM v6 made 1070 calls (100\% valid JSON, no invalid-target refusal, 0.51\,s mean latency). The median speed of the
vehicles it targeted unnecessarily was 16.0\,m/s, i.e.\ mostly ordinary traffic. In stall\_exit, 20 of its 27 successes
came from episodes in which it never targeted the stalled car: it ``bypassed'' a normal car and the resulting lane change
avoided the stalled car by chance (9 of 31 in stall\_mid). Because a maneuver runs until its target is behind the ego,
a wrong target also keeps the model from being queried again: in slow\_exit, the last decision of 38 of 50 episodes was
a maneuver on a non-injected vehicle (29 of them collided).

\begin{table}[h]
\centering\footnotesize
\caption{Paired comparisons on the same scenarios and seeds (discordant pairs; exact McNemar test).}
\label{tab:paired}
\begin{tabular}{llrrr}
\toprule
Comparison, scenario group & outcome & v6 better & other better & $p$ \\
\midrule
v6 vs DT, need & success & 58 & 0 & $<10^{-17}$ \\
v6 vs DT, no need & success & 4 & 6 & 0.75 \\
v6 vs naive, need & success & 1 & 11 & 0.006 \\
v6 vs naive, need & no collision & 4 & 20 & 0.002 \\
v6 vs v5, need & no collision & 2 & 54 & $<10^{-13}$ \\
v6 vs v5, no need & success & 11 & 1 & 0.006 \\
v6 vs designed, need & no collision & 5 & 52 & $<10^{-10}$ \\
\bottomrule
\end{tabular}
\end{table}

\paragraph{Supplementary check with larger models (offline, the same 74 development states).} Qwen3-30B-A3B-Instruct:
29.7\% correct (4B: 48.6\%). Where intervention is needed both reach 44\%. Where it is not, the 30B model intervened in
81\% of the states, and in stall\_past\_exit and stall\_upstream it targeted the stopped car in 20 of 20 states: it
recognises the abnormal object but does not use its relation to the route. Qwen3-8B (no thinking) was not better than
the 4B model; Qwen3-4B-Thinking needed about 30\,s per call and exhausted a 4096-token budget before answering, which
rules it out for a 2\,Hz closed loop on this GPU.

\paragraph{A real run (test set, stall\_mid, seed 602).} At $t=1$ (0.5\,s) the DT proposes \texttt{IDLE}. The input
and the response below are the logged ones; the executor's steps are reproduced by replaying the logged decisions,
and the replay matches the logged episode.
\end{revblock}

\begin{ModelInput}[formatcom=\color{red}]
EGO: on the south access road, 20 m before the point where it merges into the ring's outer lane,
    speed 12.5 m/s.
ROUTE: merge into the ring's outer lane; your exit, the north exit road, leaves the ring 55 m
    past your merge point.
POLICY PROPOSES: IDLE (current target speed 16 m/s; target speed levels 0 / 8 / 16 m/s; SLOWER
    and FASTER change the level, LANE_LEFT = towards the inner lane).
PERCEPTION, vehicles within 50 m. Ring positions are along the ring in the driving direction,
    relative to your merge point (+ = ahead/downstream, - = behind/upstream):
- Your road (south access road), + = ahead of you: none
- Ring, outer lane (the lane you merge into): id 3 at -31 m, 16.0 m/s; [YOUR MERGE POINT at 0
    m]; id 2 at +15 m, 16.2 m/s; id 5 at +28 m, 0.0 m/s; [YOUR EXIT (north exit road) at +55 m]
- Ring, inner lane: id 4 at -46 m, 16.2 m/s; id 1 at +8 m, 16.0 m/s; [YOUR EXIT (north exit
    road) at +55 m]
- Other roads: none
\end{ModelInput}

\begin{ModelResponse}[formatcom=\color{red}]
{"consequence": "Policy's idle action risks collision with id 5 at 0.0 m/s in outer lane; gap
    unsafe.", "decision": "BYPASS", "target_id": "5"}
\end{ModelResponse}

\begin{revblock}
\noindent Id~5 is the injected stalled car (reference: \texttt{BYPASS}, target~5). The executor accepts the id and
starts the bypass. The DT merges ($t=1$--$2$). At $t=3$ the ego is in the outer lane with the stalled car 32\,m ahead.
At $t=4$ the gap check passes and \texttt{LANE\_LEFT} is executed while the target speed is lowered to 12.1\,m/s. From
$t=5$ to $7$ the ego passes the car in the inner lane at 11.7--13.1\,m/s. At $t=9$ the car is 10\,m behind, the maneuver
ends and control returns to the DT. At $t=11$ the model answers \texttt{FOLLOW\_POLICY}; the ego exits and arrives
after 7.5\,s without collision. This run is a success and is not typical of v6; Appendix D shows a
typical failure.

\subsection*{2.6 What the results support}
\begin{enumerate}[nosep]
\item \textbf{Supported:} the integrated system (frozen DT + 4B LLM + executor) is far better than the DT where
  intervention is needed (success 0 $\rightarrow$ 38.7\%, collisions 100 $\rightarrow$ 46\%), and it is not
  significantly different from the DT in nominal traffic (96 vs.\ 94\%) and where no intervention is needed (92 vs.\ 94\%).
\item \textbf{Not supported:} that this gain comes from the model's understanding. A context-free trigger with the
  same executor is significantly better where intervention is needed (success 45.3\%, collisions 35.3\%; $p=0.006$ and
  $0.002$). v6 chose the correct decision and target in only 8\% (stalled car before the exit) and 0\% (crawler before
  the exit) of the moments that required it, started a maneuver in every episode including nominal traffic, and most
  of its successes came from incidental lane changes.
\item The low collision rate of the previous version came from the labels and the rule prompt. With them, v5 bypassed
  the stalled car \emph{after} the exit in 48 of 50 episodes (success 64\%), because its path label did not stop at the
  exit: the judgment sat in the labels, not in the model.
\item The value that contextual understanding could add in this testbed is the gap between the naive and the designed
  rule (slow\_exit collisions 80 $\rightarrow$ 18\%; stall\_past\_exit success 80 $\rightarrow$ 88\%). LLM v6 recovered none
  of it (slow\_exit collisions 80\%, the same as the naive trigger).
\end{enumerate}
With the models that run on the available hardware (4B, 8B, 30B-A3B) and a structured, facts-only input, the evidence
supports the integration effect only; it does \emph{not} show that the LLM's pretrained understanding helps the DT handle
OOD situations. \textbf{Minimal supplementary validation:} evaluate a stronger model offline on the frozen single-step
states (the 74 development states plus the states in the test decision logs). Re-run the closed loop for the
\texttt{llm} condition only if the model is clearly better than the naive trigger where intervention is needed and
close to the designed rule where it is not; all other conditions can be reused unchanged.

\paragraph{Limitations.} The input is structured (lane association, along-road position and speed come from perception);
no VLM was tested. A stalled car in the other ring lane could not serve as a clean control, because simulated traffic
crashes into it and the wreck slides into the ego's lane. Maneuvers are not re-queried while active, so a wrong target
cannot be cancelled; this applies to all methods but hurts a supervisor that often picks wrong targets most.
50 episodes per cell resolve only large differences; the conclusions use the significant paired comparisons.
\end{revblock}

\appendix
\section*{Appendix}

\subsection*{\rev{A. Scenario definitions}}
\begin{revblock}
\begin{table}[h]
\centering\footnotesize
\begin{tabular}{l>{\raggedright\arraybackslash}p{7.2cm}ll}
\toprule
Scenario & Injected vehicle (lane segment, position) & Intervention needed & DT (test) \\
\midrule
nominal & none & no & 94\% success \\
stall\_mid & outer lane (ex, ee), $s=10$\,m, 27.6\,m after the merge point, 0\,m/s & bypass & 100\% collision \\
stall\_exit & outer lane (ee, nx), $s=8$\,m, 9.6\,m before the exit, 0\,m/s & bypass & 100\% collision \\
slow\_exit & same place, 3\,m/s, follows its lane into the exit & follow & 100\% collision \\
stall\_past\_exit & outer lane (nx, ne), $s=9.6$\,m, 9.6\,m after the exit, 0\,m/s & no & 88\% success \\
stall\_upstream & outer lane (sx, se), $s=8$\,m, 12\,m upstream of the merge point, 0\,m/s & no & 100\% success \\
\bottomrule
\end{tabular}
\end{table}
Traffic within 8\,m of the injected vehicle is removed at reset. Tried and excluded: a stalled car in the inner lane
(becomes a wreck in the ego's lane) and a stalled car 12\,m after the merge point (every method waits at the entry and
times out).
\end{revblock}

\subsection*{\rev{B. System prompt (frozen, \texttt{SYSTEM\_V6\_NOCO})}}
\begin{PromptText}[formatcom=\color{red}]
You supervise a learned driving policy that controls a car (the EGO) in a two-lane roundabout.
The policy was trained by imitation only on ordinary flowing traffic in this roundabout. It drives
    well in such traffic
(merging, keeping gaps, following other cars), but it may act wrongly in situations it never saw
    during training.
At every step you decide whether the policy keeps driving or one of two maneuvers takes over.

ROAD
- The ring has two lanes (inner and outer). Traffic on the ring drives counter-clockwise. Both ring
    lanes connect to every exit.
- Four access roads lead into the ring and four exit roads lead out of it (south, east, north,
    west). Each of these roads
  has a single lane; lane changes are only possible on the ring.
- Speed limit 16 m/s.
- The ego's route: south access road -> merge into the ring's outer lane -> counter-clockwise along
    the ring -> leave on
  the north exit road. Once it has left the ring it does not come back.

INPUT (every 0.5 s): the ego's state, the rest of its route, the policy's proposed action, and every
    vehicle within
50 m detected by perception, grouped by road and lane (id, position along the road, speed).

DECISIONS
- FOLLOW_POLICY: the policy's proposed action is executed.
- BYPASS with target_id: pass the target vehicle using the other ring lane. The executor changes
    lane when the gap in that
  lane is safe, drives past the target and hands control back to the policy once the target is
    behind the ego.
  While the ego is still on the access road, the policy merges first and the lane change happens on
    the ring.
- FOLLOW_VEHICLE with target_id: stay behind the target vehicle at a safe distance without changing
    lane, slowing down
  or stopping as needed. Control returns to the policy when the target is no longer ahead of the
    ego.
The maneuvers handle braking and gap checks. They do not judge the traffic situation; that is your
    job.
Every intervention has a cost: an unnecessary lane change or slowdown delays the ego and can create
    conflicts with
other traffic.

YOUR TASK: judge the situation for the ego's own route, in this order:
1. consequence: what will happen if the policy's proposed action is executed? Will something block
    the ego or will it
   collide with something?
2. decision and target_id: the decision that lets the ego reach its exit safely and without
    unnecessary delay.

Return only JSON (each text field at most 25 words):
{"consequence": "...", "decision": "FOLLOW_POLICY | BYPASS | FOLLOW_VEHICLE", "target_id": <vehicle
    id or null>}
\end{PromptText}

\subsection*{\rev{C. Development record}}
\begin{revblock}
Offline single-step accuracy on development states (balanced per scenario and reference label), Qwen3-4B unless noted:
flat vehicle list sorted by distance 41\%; lane-grouped list with an extra ``critical objects'' field 8.7\% (the model
intervened almost everywhere); plus kinematics without that field 53\% (Qwen3-8B 47\%); plus map landmarks and a
reference no longer cut at a fixed distance 48.6\% (frozen). No further tuning was done, so that rules or answers would
not be re-introduced through the prompt. All drafts, their outputs and the reasons for each change are logged in
\texttt{reports/iteration\_log.md} (Phase~5).
\end{revblock}

\subsection*{\rev{D. A typical failure (test set, slow\_exit, seed 600, $t=1$)}}
\begin{ModelInput}[formatcom=\color{red}]
EGO: on the south access road, 20 m before the point where it merges into the ring's outer lane,
    speed 12.5 m/s.
ROUTE: merge into the ring's outer lane; your exit, the north exit road, leaves the ring 55 m
    past your merge point.
POLICY PROPOSES: SLOWER (current target speed 16 m/s; target speed levels 0 / 8 / 16 m/s; SLOWER
    and FASTER change the level, LANE_LEFT = towards the inner lane).
PERCEPTION, vehicles within 50 m. Ring positions are along the ring in the driving direction,
    relative to your merge point (+ = ahead/downstream, - = behind/upstream):
- Your road (south access road), + = ahead of you: none
- Ring, outer lane (the lane you merge into): id 6 at -30 m, 13.3 m/s; id 4 at -22 m, 16.1 m/s;
    [YOUR MERGE POINT at 0 m]; id 3 at +14 m, 16.1 m/s; [YOUR EXIT (north exit road) at +55 m]
- Ring, inner lane: id 7 at -49 m, 15.9 m/s; id 5 at -21 m, 13.5 m/s; id 2 at -11 m, 13.6 m/s;
    id 1 at +7 m, 15.7 m/s; [YOUR EXIT (north exit road) at +55 m]
- Other roads: none
\end{ModelInput}

\begin{ModelResponse}[formatcom=\color{red}]
{"consequence": "Policy slows unnecessarily; traffic flow blocked; delay and risk of collision
    with inner lane vehicles.", "decision": "BYPASS", "target_id": "4"}
\end{ModelResponse}

\begin{revblock}
\noindent The crawling car is not yet within the 50\,m sensing range. Id~4 is an ordinary car 22\,m upstream of the merge
point. The executor accepts the id and starts the bypass. The maneuver stays active until the end of the episode, so the
model is not queried again and never sees the crawler. The executor brakes for its lane leader but finds no gap for the
lane change. At $t=15$ (7.5\,s), at 3.2\,m/s and about 10\,m behind the crawler, the ego has an at-fault collision. On the
same seed neither the naive trigger nor the designed rule collides.
\end{revblock}

\subsection*{\rev{E. Code and data}}
\begin{revblock}
Repository \url{https://github.com/Zhang1998-06/ce-llm-roundabout}. Scenarios: \texttt{src/hazards.py}; facts-only input,
prompt and query: \texttt{src/coc\_supervisor.py} (\texttt{perceive\_v6}, \texttt{describe\_v6}, \texttt{SYSTEM\_V6\_NOCO},
\texttt{ask\_v6}); controls and reference: \texttt{src/rules\_v6.py}; closed loop: \texttt{scripts/run\_hazard.py};
tables: \texttt{scripts/make\_tables\_v6.py}; replay: \texttt{scripts/replay\_v6.py}; data: \texttt{results/hazard\_test\_v6/}
(per-episode outcomes and every supervisor call with the full user message and raw response).
\end{revblock}

\end{document}
